\section{Inputs for the threshold construction}
\label{sec:rhost-threshold-tools}

All graphs are simple.
A path's length is its number of edges.
The neighbourhood $N_J(A)$ is external: it contains vertices outside $A$ adjacent to $A$.
The cycle rank of a connected graph is $|E|-|V|+1$.
We write $\ln$ for the natural logarithm and $\log_2$ for the binary logarithm.
An event holds with high probability if its probability tends to one as $n\to\infty$.

\begin{lemma}[Uniform robust expansion]\label{lem:rhost-threshold-expansion}
Suppose $D\to\infty$, $D\le(\ln n)^B$ for a fixed $B$, and $\beta=20\ln D/D$.
Let $R\subseteq[n]$ have size $n-o(n)$ and expose $J\sim G(R,D/n)$.
For any prescribed $r_n=o(n)$, with high probability every $U\subseteq R$ of size at least $n-r_n$
contains a set $U'$ with
\[
 |U'|>|U|-\beta n,\qquad
 |N_{J[U']}(A)|\ge\frac{|A|}{3\beta}
 \quad(0<|A|\le\beta n).
\]
The assertion also holds conditionally on a random $R$ chosen independently of $J$.
\end{lemma}

\begin{proof}
Put $s=\lceil\beta n\rceil$.
The probability that two disjoint $s$-sets in $R$ have no joining edge is at most
\[
 \binom ns^2e^{-Ds^2/n}
 \le\exp\{s[2\ln(en/s)-Ds/n]\}=o(1).
\]
Indeed, $Ds/n\sim20\ln D$ whereas $\ln(en/s)\sim\ln D$.
Assume no such pair exists.
Starting from $U$, repeatedly remove a nonempty set of at most $\beta n$ vertices that violates the desired expansion.
If the total removed set $B$ first reaches size $\beta n$, then $\beta n\le|B|<2\beta n$.
Its external neighbourhood in the original $J[U]$ has size less than $|B|/(3\beta)<2n/3$.
Thus $U\setminus(B\cup N_{J[U]}(B))$ has at least $n/3-r_n-2\beta n$ vertices.
For large $n$, this exceeds $s$.
Taking $s$ vertices there and $s$ vertices in $B$ gives the forbidden pair.
Hence the deletion process removes fewer than $\beta n$ vertices and proves the assertion.
The probability estimate is uniform in $R$.
\end{proof}

The rooted tree theorem below is the form of Haxell's embedding result used in the construction.
It is stated as Theorem~2.2 by Alon and Krivelevich,
\href{https://doi.org/10.1137/23M1598969}{\emph{Sparse Pancyclic Subgraphs of Random Graphs}} (2025).
That statement follows from Haxell, \emph{Tree embeddings}, Journal of Graph Theory 36 (2001), 121--130,
and Balogh, Csaba, Pei, and Samotij,
\href{https://doi.org/10.37236/278}{\emph{Large bounded degree trees in expanding graphs}} (2010).

\begin{theorem}[Rooted tree embedding, external]\label{thm:rhost-rooted-tree}
Let $N,\Delta,k$ be positive integers.
Suppose a graph $J$ satisfies
\[
 |N_J(A)|\ge\Delta|A|+1\quad(0<|A|\le k),
\]
\[
 |N_J(A)|\ge\Delta|A|+N\quad(k<|A|\le2k).
\]
Then every tree with at most $N$ vertices and maximum degree at most $\Delta$
embeds in $J$ with any specified root mapped to any specified vertex of $J$.
\end{theorem}

\begin{lemma}[Trees with many leaves at one depth]\label{lem:rhost-threshold-trees}
Use $\beta$ from Lemma~\ref{lem:rhost-threshold-expansion} and set
\[
 a=\lfloor(5\beta)^{-1}\rfloor,\quad M=\lfloor n/1000\rfloor,
 \quad d=\lceil\log_a M\rceil.
\]
Every expander $J[U']$ supplied by that lemma contains, rooted at any prescribed vertex,
a tree with exactly $M$ leaves at depth $d$, fewer than $3M$ vertices, and maximum degree at most $a+1$.
It also contains a tree formed by joining the roots of two such trees with any positive-length $o(n)$ stem.
\end{lemma}

\begin{proof}
Take $k=\lfloor\beta n/2\rfloor$, $\Delta=\lfloor1/(4\beta)\rfloor$, and $N=\lfloor n/50\rfloor$.
For $0<|A|\le k$, the expansion surplus over $\Delta|A|$ is at least $|A|/(12\beta)>1$.
For $k<|A|\le2k$, it exceeds $n/24>N$.
Thus Theorem~\ref{thm:rhost-rooted-tree} applies.
For sufficiently large $n$, we have $a\ge3$ and $a+1\le\Delta$.

Build complete $a$-ary levels through depth $d-1$.
Give the $a^{d-1}$ vertices at that depth between one and $a$ children each, with exactly $M$ children in total.
This is possible because $a^{d-1}\le M\le a^d$.
All leaves have depth $d$.
The earlier levels have fewer than $2M$ vertices, so the whole tree has fewer than $3M$ vertices.
Two such trees, joined at their roots by a positive-length $o(n)$ stem, have fewer than $6M+o(n)<N$ vertices.
Their maximum degree is at most $a+1$.
Apply the rooted embedding theorem to each required tree.
\end{proof}

\begin{lemma}[A short list of disjoint cycles]\label{lem:rhost-cycle-list}
Fix constants $A,B>0$.
Suppose $D\to\infty$ and $D\le(\ln n)^B$.
In $G(n,D/n)$, with high probability, one can select vertex-disjoint cycles from any fixed list of $O(\ln n)$ lengths,
all between $3$ and $A\ln n$, provided each length has bounded multiplicity.
\end{lemma}

\begin{proof}
Let $X_\ell$ count cycles of length $\ell=O(\ln n)$ and let $Y_\ell$ count ordered intersecting pairs, including identical copies.
Write $\mu_\ell=\mathbb EX_\ell$.
We first establish
\[
 \operatorname{Var}X_\ell,\ \mathbb EY_\ell
 \le C\mu_\ell^2\left(\frac{\ell^2}n+\ell D^{-\ell}\right).
\]
Every proper overlap is a union of paths and isolated vertices.
An overlap with $u$ vertices and $f$ edges has normalized weight $(1+O(\ell^2/n))n^{-u}(n/D)^f$.
An isolated vertex has at most $\ell^2$ descriptions.
A path with $j\ge1$ edges has at most $4\ell^2$ descriptions, each of weight $n^{-1}D^{-j}$.
For several components, division by their order factorial bounds the total by the exponential of the one-component sum.
Thus all proper nonempty overlaps contribute $O(\ell^2/n)$.
The full overlap contributes $(1+O(\ell^2/n))2\ell D^{-\ell}$.
Vertex-disjoint pairs contribute at most $\mu_\ell^2$, proving the variance bound as well.

Put $q=\lfloor\ln n/(10\ln D)\rfloor$.
Chebyshev's inequality, summed over $3\le\ell\le q$, shows that every required short length has any fixed number of copies.
Indeed, $\sum_{\ell\ge3}\ell D^{-\ell}=O(D^{-3})$, and $\mu_\ell\to\infty$ uniformly.
With high probability, all distinct cycles of lengths at most $q$ are vertex-disjoint.
The same proper-overlap enumeration for two possibly different lengths bounds the expected number of intersecting pairs by
\[
 n^{-1}(\ln n)^{O(1)}D^{2q}=o(1).
\]
Choose all prescribed cycles in this range first.

For $q<\ell\le A\ln n$, put $\varepsilon_\ell=C(\ell^2/n+\ell D^{-\ell})\le n^{-1/10+o(1)}$.
Let $M_0=\lceil(\ln n)^3\rceil$.
Chebyshev's inequality gives $X_\ell\ge\mu_\ell/2$ except with probability $O(\varepsilon_\ell)$.
Markov's inequality gives $Y_\ell\le\mu_\ell^2/(8M_0)$ except with probability $O(M_0\varepsilon_\ell)$.
In the graph whose vertices are copies and whose edges join intersecting copies, the elementary independent-set bound is $X_\ell^2/Y_\ell$.
Thus at least $2M_0$ copies are pairwise vertex-disjoint.
A union bound gives this simultaneously for all required lengths.
The prescribed list uses only $O((\ln n)^2)$ vertices in total.
Greedily choose the remaining cycles: each previously used vertex excludes at most one member of any disjoint-copy family.
\end{proof}

The Hamilton-path input below is Lemma~2.1 in the cited Alon--Krivelevich paper.
It follows from the Hamiltonicity threshold and rotation-extension arguments; see Section~4 of
Krivelevich, \href{https://arxiv.org/abs/1507.00205}{\emph{Long paths and Hamiltonicity in random graphs}} (2016).

\begin{lemma}[Hamilton completion, external input]\label{lem:rhost-threshold-hamilton}
If $Nq-\ln N-\ln\ln N\to\infty$, then $G(N,q)$ has, with high probability,
a set $S$ of at least $N/4$ vertices such that every $s\in S$ has a set $T_s$ of at least $N/4$ vertices,
and each $t\in T_s$ is joined to $s$ by a Hamilton path.
Consequently, two external vertices with fresh independent $q$-edges to this graph can be joined through all $N$ vertices with high probability.
\end{lemma}

\begin{proof}[Proof of the consequence]
Expose the induced graph first.
The first external vertex has a neighbour $s\in S$ except with probability at most $e^{-qN/4}$.
Condition on that neighbour.
The second external vertex has a neighbour in $T_s$ with the same failure bound.
These two edges and the stated Hamilton path give the required path.
\end{proof}

\begin{lemma}[Long interval chords]\label{lem:rhost-threshold-windows}
Let a Hamilton cycle on $n$ vertices contain a distinguished arc of length
$2s-1+\ell-1$, where $s\gg n^{0.9}$ and $\ell=o(n^{0.9})$.
Suppose replacing paths inside this arc permits every length reduction from $0$ through $2s-1$.
Let $m=\lfloor n/s\rfloor\ge3$.
In an independent $G(n,D/n)$ with $D\to\infty$, with high probability $m-2$ additional edges
supply cycle intervals that, together with
\[
 [\ell,\ell+2s-1]\quad\text{and}\quad[n-2s+1,n],
\]
cover every integer from $\ell$ through $n$.
\end{lemma}

\begin{proof}
Write $W=n^{0.9}$ and let the distinguished arc have length $b_0=2s+\ell-2$.
For each $3\le i\le m$, there are at least $W^2/100$ vertex pairs whose joining arc contains the distinguished arc
and has length $a_i\in[is-W,is+W]$.
To see this, extend the two ends of the distinguished arc by distances $u,v\ge0$.
Put $t_i=\min\{is-b_0,n-b_0-W/2\}$.
Vary each distance independently in a window of width $W/4$ centered at $t_i/2$.
The lower bound $is-b_0\ge s-o(s)\gg W$ keeps both distances positive.
The upper adjustment leaves a complementary arc of length at least $W/4$.
All these pairs are distinct non-cycle edges and satisfy the required length window.

The probability that some window has no random edge is at most
\[
 m\exp(-Dn^{0.8}/100)=o(1).
\]
Choose one edge from each window.
It closes a cycle on the arc of length $a_i$, and the allowed reductions give
$[a_i+1-(2s-1),a_i+1]$.
Consecutive intervals overlap because $s\gg W$.
The first overlaps $[\ell,\ell+2s-1]$, and the last overlaps $[n-2s+1,n]$ because $ms>n-s$.
This proves the assertion.
\end{proof}

Finally, every connected pancyclic graph with $n+k$ edges has cycle rank $k+1$.
Its distinct cycles give distinct nonzero vectors in its binary cycle space.
Hence $n-2\le2^{k+1}-1$, so $k\ge\log_2(n-1)-1$.
